%! AP Statistics — Unit 3, Lesson 3.2 — Homework KEY
\documentclass[10pt]{article}
\usepackage{apstats-article}
\usepackage{apstats-key}

\begin{document}

\pageheader{Unit 3, Lesson 3.2}{Homework — KEY}
\namedateperiod

\vspace{0.1in}

\begin{practicebox}
\small
\textbf{1.} Diet and blood pressure.

Study type: \ans{Observational study.} \quad Rand. selection: \ans{Yes.} \quad Rand. assignment: \ans{No.}\\[4pt]
Causation: \ans{No.} \quad Generalize: \ans{Yes (to city adults).}\\[4pt]
\ans{This is an observational study — the researcher observed but did not impose diets. Confounding variables (exercise, genetics) could explain both diet quality and blood pressure. Random selection allows generalization to the city adult population.}

\vspace{6pt}
\textbf{2.} Standing desk experiment.

Study type: \ans{Experiment.} \quad Rand. selection: \ans{No (recruited volunteers).} \quad Rand. assignment: \ans{Yes.}\\[4pt]
Causation: \ans{Yes.} \quad Generalize: \ans{No (only volunteers at this company).}\\[4pt]
\ans{Random assignment was used, so the standing desk caused the reduction in back pain. However, employees were not randomly selected, so results cannot generalize beyond similar employees at this company.}

\vspace{6pt}
\textbf{3.} Reading intervention.
\begin{enumerate}[label=(\alph*), leftmargin=2em, itemsep=5pt]
  \item \ans{Yes — random assignment was used, so the intervention caused the significantly higher scores.}
  \item \ans{No — students were volunteers from a single school, not randomly selected from all third-graders. Results cannot generalize beyond similar volunteers.}
  \item \ans{Randomly select students from a broader population of third-graders (e.g., from a district registry) rather than recruiting volunteers from one school.}
\end{enumerate}
\end{practicebox}

\end{document}
